\documentclass{ximera}

\begin{document}
	\author{Stitz-Zeager}
	\xmtitle{Exercises for Trigonometric Equations and Inequalities}{}

\mfpicnumber{1} \opengraphsfile{ExercisesforTrigonometricEquationsandInequalities} % mfpic settings added 

\begin{question}
In Exercises \ref{solvebasicfirst} - \ref{solvebasiclast}, find \underline{all} of the exact solutions of the  equation and then list those solutions which are in the interval $[0, 2\pi)$.

\begin{problem} \label{solvebasicfirst}
$\sin \left( 5 \theta \right) = 0$ 

\begin{solution}
$\theta = \frac{\pi k}{5}; \; \theta = 0, \frac{\pi}{5}, \frac{2\pi}{5}, \frac{3\pi}{5}, \frac{4\pi}{5}, \pi, \frac{6\pi}{5}, \frac{7\pi}{5}, \frac{8\pi}{5}, \frac{9\pi}{5}$
\end{solution}
\end{problem}

\begin{problem}
$\cos \left( 3t \right) = \frac{1}{2}$

\begin{solution}
$t  = \frac{\pi}{9} + \frac{2\pi k}{3}$ or $t = \frac{5\pi}{9} + \frac{2\pi k}{3}; \; t = \frac{\pi}{9}, \frac{5\pi}{9}, \frac{7\pi}{9}, \frac{11\pi}{9}, \frac{13\pi}{9}, \frac{17\pi}{9}$
\end{solution}
\end{problem}

\begin{problem}
$\sin \left( -2x \right) = \frac{\sqrt{3}}{2}$

\begin{solution}
$x = \frac{2\pi}{3} + \pi k$ or $x = \frac{5\pi}{6} + \pi k; \; x = \frac{2\pi}{3}, \frac{5\pi}{6}, \frac{5\pi}{3}, \frac{11\pi}{6}$
\end{solution}
\end{problem}

\begin{problem}
$\tan \left( 6 \theta \right) = 1$

\begin{solution}
$\theta = \frac{\pi}{24} + \frac{\pi k}{6}; \; \theta = \frac{\pi}{24}, \frac{5\pi}{24}, \frac{3\pi}{8}, \frac{13\pi}{24}, \frac{17\pi}{24}, \frac{7\pi}{8}, \frac{25\pi}{24}, \frac{29\pi}{24}, \frac{11\pi}{8}, \frac{37\pi}{24}, \frac{41\pi}{24}, \frac{15\pi}{8}$
\end{solution}
\end{problem}

\begin{problem}
$\csc \left( 4 t \right) = -1$

\begin{solution}
$t = \frac{3\pi}{8} + \frac{\pi k}{2}; \; t = \frac{3\pi}{8}, \frac{7\pi}{8}, \frac{11\pi}{8}, \frac{15\pi}{8}$
\end{solution}
\end{problem}

\begin{problem}
$\sec \left( 3x \right) = \sqrt{2}$

\begin{solution}
$x = \frac{\pi}{12} + \frac{2\pi k}{3}$ or $x = \frac{7\pi}{12} + \frac{2\pi k}{3}; \; x = \frac{\pi}{12}, \frac{7\pi}{12}, \frac{3\pi}{4}, \frac{5\pi}{4}, \frac{17\pi}{12}, \frac{23\pi}{12}$
\end{solution}
\end{problem}

\begin{problem}
$\cot \left( 2 \theta \right) = -\frac{\sqrt{3}}{3}$

\begin{solution}
$\theta  = \frac{\pi}{3} + \frac{\pi k}{2}; \; \theta  = \frac{\pi}{3}, \frac{5\pi}{6}, \frac{4\pi}{3}, \frac{11\pi}{6}$
\end{solution}
\end{problem}

\begin{problem}
$\cos \left( 9t  \right) = 9$

\begin{solution}
No solution
\end{solution}
\end{problem}

\begin{problem}
$\sin \left( \frac{x}{3} \right) = \frac{\sqrt{2}}{2}$

\begin{solution}
$x = \frac{3\pi}{4} + 6\pi k$ or $x = \frac{9\pi}{4} + 6\pi k; \; x = \frac{3\pi}{4}$
\end{solution}
\end{problem}

\begin{problem}
$\cos \left( \theta+ \frac{5\pi}{6} \right) = 0$

\begin{solution}
$\theta = -\frac{\pi}{3} + \pi k; \; \theta  = \frac{2\pi}{3}, \frac{5\pi}{3}$
\end{solution}
\end{problem}

\begin{problem}
$\sin \left( 2t - \frac{\pi}{3} \right) = -\frac{1}{2}$

\begin{solution}
$t = \frac{3\pi}{4} + \pi k$ or $t = \frac{13\pi}{12} + \pi k; \; t = \frac{\pi}{12}, \frac{3\pi}{4}, \frac{13\pi}{12}, \frac{7\pi}{4}$
\end{solution}
\end{problem}

\begin{problem}
$2\cos \left( x + \frac{7\pi}{4} \right) = \sqrt{3}$

\begin{solution}
$x = -\frac{19\pi}{12} + 2\pi k$ or $x = \frac{\pi}{12} + 2\pi k; \; x = \frac{\pi}{12}, \frac{5\pi}{12}$
\end{solution}
\end{problem}

\begin{problem}
$\csc( \theta) = 0$

\begin{solution}
No solution
\end{solution}
\end{problem}

\begin{problem}
$\tan \left( 2t - \pi \right) = 1$

\begin{solution}
$t = \frac{5\pi}{8} + \frac{\pi k}{2}; \; t = \frac{\pi}{8}, \frac{5\pi}{8}, \frac{9\pi}{8}, \frac{13\pi}{8}$
\end{solution}
\end{problem}

\begin{problem}
$\tan^{2} \left( x \right) = 3$

\begin{solution}
$x = \frac{\pi}{3} + \pi k$ or $x = \frac{2\pi}{3} + \pi k; \; x = \frac{\pi}{3}, \frac{2\pi}{3}, \frac{4\pi}{3}, \frac{5\pi}{3}$
\end{solution}
\end{problem}

\begin{problem}
$\sec^{2} \left( \theta \right) = \frac{4}{3}$

\begin{solution}
$\theta  = \frac{\pi}{6} + \pi k$ or $\theta = \frac{5\pi}{6} + \pi k; \; \theta = \frac{\pi}{6}, \frac{5\pi}{6}, \frac{7\pi}{6}, \frac{11\pi}{6}$
\end{solution}
\end{problem}

\begin{problem}
$\cos^{2} \left( t \right) = \frac{1}{2}$

\begin{solution}
$t = \frac{\pi}{4} + \frac{\pi k}{2}; \; t = \frac{\pi}{4}, \frac{3\pi}{4}, \frac{5\pi}{4}, \frac{7\pi}{4}$
\end{solution}
\end{problem}

\begin{problem} \label{solvebasiclast}
$\sin^{2} \left( x \right) = \frac{3}{4}$ 

\begin{solution}
$x = \frac{\pi}{3} + \pi k$ or $x = \frac{2\pi}{3} + \pi k; \; x = \frac{\pi}{3}, \frac{2\pi}{3}, \frac{4\pi}{3}, \frac{5\pi}{3}$
\end{solution}
\end{problem}
\end{question}

\begin{question}
In Exercises \ref{solveidentfirst} - \ref{solveidentlast}, solve the equation, giving the exact solutions which lie in $[0, 2\pi)$

\begin{problem} \label{solveidentfirst}
$\sin \left( \theta \right) = \cos \left( \theta \right)$ 

\begin{solution}
$\theta = \frac{\pi}{4}, \frac{5\pi}{4}$
\end{solution}
\end{problem}

\begin{problem}
$\sin \left( 2t \right) = \sin \left( t \right)$

\begin{solution}
$t = 0, \frac{\pi}{3}, \pi, \frac{5\pi}{3}$
\end{solution}
\end{problem}

\begin{problem}
$\sin \left( 2x \right) = \cos \left( x \right)$

\begin{selectAll}
    \choice[correct]{$x = \frac{\pi}{6}$}
    \choice{$x = \frac{\pi}{3}$}
    \choice[correct]{$x = \frac{\pi}{2}$}
    \choice{$x = \frac{2\pi}{3}$}
    \choice[correct]{$x = \frac{5\pi}{6}$}
    \choice{$x = \pi$}
    \choice[correct]{$x= \frac{3\pi}{2}$}
  \end{selectAll}

\end{problem}

\begin{problem}
$\cos \left( 2\theta \right) = \sin \left( \theta \right)$

\begin{solution}
$\theta = \frac{\pi}{6}, \frac{5\pi}{6}, \frac{3\pi}{2}$
\end{solution}
\end{problem}

\begin{problem}
$\cos \left( 2t \right) = \cos \left( t \right)$

\begin{selectAll}
    \choice[correct]{$t=0$}
    \choice{$t = \frac{\pi}{3}$}
    \choice[correct]{$t = \frac{2\pi}{3}$}
    \choice{$t = \pi$}
    \choice[correct]{$t = \frac{4\pi}{3}$}
    \choice{$t = \frac{5\pi}{3}$}
  \end{selectAll}

\end{problem}

\begin{problem}
$\cos(2x) = 2 - 5\cos(x)$

\begin{solution}
$x=\frac{\pi}{3}, \frac{5\pi}{3}$
\end{solution}
\end{problem}

\begin{problem}
$3\cos(2 \theta ) + \cos(\theta) + 2 = 0$

\begin{solution}
$\theta = \frac{2\pi}{3}, \frac{4\pi}{3}, \arccos\left(\frac{1}{3}\right), 2\pi -\arccos\left(\frac{1}{3}\right) $
\end{solution}
\end{problem}

\begin{problem}
$\cos(2t) = 5\sin(t) - 2$

\begin{selectAll}
    \choice[correct]{$t = \frac{\pi}{6}$}
    \choice{$t = \frac{\pi}{3}$}
    \choice{$t = \frac{\pi}{2}$}
    \choice{$t = \frac{2\pi}{3}$}
    \choice[correct]{$t = \frac{5\pi}{6}$}
    \choice{$t = \pi$}
    \choice{$t = \frac{3\pi}{2}$}
  \end{selectAll}
\end{problem}

\begin{problem}
$3\cos(2x) = \sin(x) + 2$

\begin{solution}
$x = \frac{7\pi}{6}, \frac{11\pi}{6}, \arcsin\left(\frac{1}{3}\right), \pi - \arcsin\left(\frac{1}{3}\right) $
\end{solution}
\end{problem}
 
\begin{problem} 
$2\sec^{2}(\theta) = 3 - \tan(\theta)$

\begin{solution}
$\theta=\frac{3\pi}{4}, \frac{7\pi}{4}, \arctan\left(\frac{1}{2}\right), \pi +\arctan\left(\frac{1}{2}\right) $
\end{solution}
\end{problem}

\begin{problem}
$\tan^{2}(t) = 1-\sec(t)$

\begin{selectAll}
    \choice[correct]{$t=0$}
    \choice{$t = \frac{\pi}{3}$}
    \choice[correct]{$t = \frac{2\pi}{3}$}
    \choice{$t = \pi$}
    \choice[correct]{$t = \frac{4\pi}{3}$}
    \choice{$t = \frac{5\pi}{3}$}
  \end{selectAll}
\end{problem}


\begin{problem}
$\cot^{2}(x) = 3\csc(x) - 3$

\begin{selectAll}
    \choice[correct]{$x = \frac{\pi}{6}$}
    \choice{$x = \frac{\pi}{3}$}
    \choice[correct]{$x = \frac{\pi}{2}$}
    \choice{$x = \frac{2\pi}{3}$}
    \choice[correct]{$x = \frac{5\pi}{6}$}
    \choice{$x = \pi$}
  \end{selectAll}
\end{problem}

\begin{problem}
$\sec(\theta) = 2\csc(\theta)$

\begin{solution}
$\theta=\arctan(2), \pi + \arctan(2)$ 
\end{solution}
\end{problem}

\begin{problem}
$\cos(t) \csc(t)\cot(t) = 6-\cot^{2}(t)$

\begin{solution}
$t = \frac{\pi}{6}, \frac{7\pi}{6}, \frac{5\pi}{6}, \frac{11\pi}{6}$
\end{solution}
\end{problem}

\begin{problem}
$\sin(2x) = \tan(x)$

\begin{solution}
$x = 0, \pi, \frac{\pi}{4}, \frac{3\pi}{4}, \frac{5\pi}{4}, \frac{7\pi}{4}$
\end{solution}
\end{problem}

\begin{problem}
$\cot^{4}(\theta) = 4\csc^{2}(\theta) - 7$

\begin{solution}
$\theta = \frac{\pi}{6}, \frac{\pi}{4}, \frac{3\pi}{4}, \frac{5\pi}{6}, \frac{7\pi}{6}, \frac{5\pi}{4}, \frac{7\pi}{4}, \frac{11\pi}{6}$
\end{solution}
\end{problem}

\begin{problem}
 $\cos(2t) + \csc^{2}(t) = 0$

\begin{solution}
$t = \frac{\pi}{2}, \frac{3\pi}{2}$
\end{solution}
\end{problem}

\begin{problem}
$\tan^{3} \left( x \right) = 3\tan \left( x \right)$

\begin{solution}
$x = 0, \frac{\pi}{3}, \frac{2\pi}{3}, \pi, \frac{4\pi}{3}, \frac{5\pi}{3}$
\end{solution}
\end{problem}

\begin{problem}
$\tan^{2} \left( \theta \right) = \frac{3}{2} \sec \left( \theta \right)$

\begin{solution}
$\theta  = \frac{\pi}{3}, \frac{5\pi}{3}$
\end{solution}
\end{problem}

\begin{problem}
$\cos^{3} \left( t \right) = -\cos \left( t \right)$ 

\begin{solution}
$t = \frac{\pi}{2}, \frac{3\pi}{2}$
\end{solution}
\end{problem}

\begin{problem}
$\tan (2x) - 2\cos(x) = 0$

\begin{solution}
$x = \frac{\pi}{6}, \frac{\pi}{2}, \frac{5\pi}{6}, \frac{3\pi}{2}$
\end{solution}
\end{problem}

\begin{problem}
$\csc^{3}(\theta) + \csc^{2}(\theta) = 4\csc(\theta) + 4$

\begin{solution}
$\theta = \frac{\pi}{6}, \frac{5\pi}{6}, \frac{7\pi}{6}, \frac{3\pi}{2}, \frac{11\pi}{6}$
\end{solution}
\end{problem}

\begin{problem}
$2\tan(t) = 1 - \tan^{2}(t)$

\begin{solution}
$t = \frac{\pi}{8}, \frac{5\pi}{8}, \frac{9\pi}{8}, \frac{13\pi}{8}$
\end{solution}
\end{problem}

\begin{problem} \label{solveidentlast}
$\tan \left( x \right) = \sec \left( x \right)$ 

\begin{solution}
No solution 
\end{solution}
\end{problem}

\end{question}

\begin{question}
In Exercises \ref{solvemoreidentfirst} - \ref{solvemoreidentlast}, solve the equation, giving the exact solutions which lie in $[0, 2\pi)$

\begin{problem} \label{solvemoreidentfirst}
$\sin(6\theta) \cos(\theta) = -\cos(6\theta) \sin(\theta)$ 

\begin{solution}
$\theta = 0, \frac{\pi}{7}, \frac{2\pi}{7}, \frac{3\pi}{7}, \frac{4\pi}{7}, \frac{5\pi}{7}, \frac{6\pi}{7}, \pi, \frac{8\pi}{7}, \frac{9\pi}{7}, \frac{10\pi}{7}, \frac{11\pi}{7}, \frac{12\pi}{7}, \frac{13\pi}{7}$
\end{solution}
\end{problem}

\begin{problem}
$\sin(3t)\cos(t) = \cos(3t) \sin(t)$

\begin{solution}
$t=0, \frac{\pi}{2}, \pi, \frac{3\pi}{2}$
\end{solution}
\end{problem}

\begin{problem}
$\cos(2x)\cos(x) + \sin(2x)\sin(x) = 1$ 

\begin{solution}
$x = 0$
\end{solution}
\end{problem}

\begin{problem}
$\cos(5\theta)\cos(3\theta) - \sin(5\theta)\sin(3\theta) = \frac{\sqrt{3}}{2}$

\begin{solution}
$\theta = \frac{\pi}{48}, \frac{11\pi}{48}, \frac{13\pi}{48}, \frac{23\pi}{48}, \frac{25\pi}{48}, \frac{35\pi}{48}, \frac{37\pi}{48}, \frac{47\pi}{48}, \frac{49\pi}{48}, \frac{59\pi}{48}, \frac{61\pi}{48}, \frac{71\pi}{48}, \frac{73\pi}{48}, \frac{83\pi}{48}, \frac{85\pi}{48}, \frac{95\pi}{48}$
\end{solution}
\end{problem}

\begin{problem}
$\sin(t) + \cos(t) = 1$

\begin{selectAll}
    \choice[correct]{$t = 0$}    
    \choice{$t = \frac{\pi}{6}$}
    \choice{$t = \frac{\pi}{3}$}
    \choice[correct]{$t = \frac{\pi}{2}$}
    \choice{$t = \frac{2\pi}{3}$}
    \choice{$t = \frac{5\pi}{6}$}
    \choice{$t = \pi$}
  \end{selectAll}
\end{problem}

\begin{problem}
$\sin(x) + \sqrt{3} \cos(x) = 1$

\begin{solution}
$x = \frac{\pi}{2}, \frac{11\pi}{6}$
\end{solution}
\end{problem}

\begin{problem}
$\sqrt{2} \cos(\theta) - \sqrt{2} \sin(\theta) = 1$

\begin{solution}
$\theta = \frac{\pi}{12}, \frac{17\pi}{12}$
\end{solution}
\end{problem}

\begin{problem}
$\sqrt{3} \sin(2t) +  \cos(2t) = 1$

\begin{selectAll}
    \choice[correct]{$t=0$}
    \choice[correct]{$t = \frac{\pi}{3}$}
    \choice{$t = \frac{2\pi}{3}$}
    \choice[correct]{$t = \pi$}
    \choice[correct]{$t = \frac{4\pi}{3}$}
    \choice{$t = \frac{5\pi}{3}$}
  \end{selectAll}
\end{problem}

\begin{problem}
 $\cos(2x) - \sqrt{3} \sin(2x) = \sqrt{2}$

\begin{solution}
$x = \frac{17 \pi}{24}, \frac{41 \pi}{24}, \frac{23\pi}{24}, \frac{47\pi}{24}$
\end{solution}
\end{problem}
 
\begin{problem}
$3\sqrt{3}\sin(3\theta) - 3\cos(3\theta) = 3\sqrt{3}$

\begin{solution}
$\theta = \frac{\pi}{6}, \frac{5\pi}{18}, \frac{5\pi}{6}, \frac{17\pi}{18}, \frac{3\pi}{2}, \frac{29\pi}{18}$
\end{solution}
\end{problem}

\begin{problem}
$\cos(3t) = \cos(5t)$

\begin{solution}
$t = 0, \frac{\pi}{4}, \frac{\pi}{2}, \frac{3\pi}{4}, \pi, \frac{5\pi}{4}, \frac{3\pi}{2}, \frac{7\pi}{4}$
\end{solution}
\end{problem}

\begin{problem}
$\cos(4x) = \cos(2x)$

\begin{solution}
$x = 0, \frac{\pi}{3}, \frac{2\pi}{3}, \pi, \frac{4\pi}{3}, \frac{5\pi}{3}$
\end{solution}
\end{problem}

\begin{problem}
$\sin(5\theta) = \sin(3\theta)$

\begin{solution}
$\theta = 0, \frac{\pi}{8}, \frac{3\pi}{8}, \frac{5\pi}{8}, \frac{7\pi}{8}, \pi, \frac{9\pi}{8}, \frac{11\pi}{8}, \frac{13\pi}{8}, \frac{15\pi}{8}$
\end{solution}
\end{problem}

\begin{problem}
$\cos(5t) = -\cos(2t)$

\begin{solution}
$t = \frac{\pi}{7}, \frac{\pi}{3}, \frac{3\pi}{7}, \frac{5\pi}{7}, \pi, \frac{9\pi}{7}, \frac{11\pi}{7}, \frac{5\pi}{3}, \frac{13\pi}{7}$ 
\end{solution}
\end{problem}

\begin{problem}
$\sin(6x) + \sin(x) = 0$

\begin{solution}
$x = 0, \frac{2\pi}{7}, \frac{4\pi}{7}, \frac{6\pi}{7}, \frac{8\pi}{7}, \frac{10\pi}{7}, \frac{12\pi}{7}, \frac{\pi}{5}, \frac{3\pi}{5}, \pi, \frac{7\pi}{5}, \frac{9\pi}{5}$
\end{solution}
\end{problem}

\begin{problem} \label{solvemoreidentlast}
$\tan(x) = \cos(x)$ 

\begin{solution}
$x = \arcsin \left( \frac{-1 + \sqrt{5}}{2} \right) \approx 0.6662, \pi - \arcsin \left( \frac{-1 + \sqrt{5}}{2} \right) \approx 2.4754$
\end{solution}
\end{problem}

\end{question}

\begin{question}
In Exercises \ref{firstinveqn} - \ref{lastinveqn}, solve the equation.

\begin{problem} \label{firstinveqn}
$\arccos(2x) = \pi$ 

\begin{solution}
$x = -\frac{1}{2}$
\end{solution}
\end{problem}

\begin{problem}
$\pi - 2\arcsin(t) = 2\pi$   

$t=\answer{-1}$
\end{problem}

\begin{problem}
$4\arctan(3x-1)-\pi=0$

\begin{solution}
$x = \frac{2}{3}$
\end{solution}
\end{problem}

\begin{problem}
$6 \, \text{arccot}(2t) - 5\pi = 0$   

$t=\answer{-\frac{\sqrt{3}}{2}}$
\end{problem}

\begin{problem}
$4 \,\text{arcsec}\left(\frac{x}{2}\right) = \pi$   

\begin{solution}
$x = 2\sqrt{2}$
\end{solution}
\end{problem}

\begin{problem}
$12 \,\text{arccsc}\left(\frac{t}{3}\right) = 2\pi$

$t = \answer{6}$
\end{problem}

\begin{problem}
$9 \arcsin^{2}(x) - \pi^2 = 0$   

\begin{solution}
$x = \pm \frac{\sqrt{3}}{2}$
\end{solution}
\end{problem}

\begin{problem}
$9 \arccos^{2}(t) - \pi^2 = 0$   

$t = \answer{\frac{1}{2}}$
\end{problem}

\begin{problem}
$8 \, \text{arccot}^{2}(x)+3\pi^2=10 \pi \, \text{arccot}(x)$   

\begin{solution}
$x = -1,0$
\end{solution}
\end{problem}
 
\begin{problem} \label{lastinveqn}
$6 \arctan(t)^2= \pi \arctan(x)+\pi^2$    

$t = \answer{-\sqrt{3}}$
\end{problem}

\end{question}

\begin{question}
In Exercises \ref{firstineqfirst} - \ref{firstineqlast}, solve the inequality.  Express the exact answer in \underline{interval} notation, restricting your attention to $0 \leq x \leq 2\pi$.

\begin{problem} \label{firstineqfirst}
$\sin \left( x \right) \leq 0$ 

\begin{solution}
$\left[ \pi, 2\pi \right]$
\end{solution}
\end{problem}

\begin{problem}
$\tan \left( t \right) \geq \sqrt{3}$

\begin{solution}
$\left[ \frac{\pi}{3}, \frac{\pi}{2} \right) \cup \left[ \frac{4\pi}{3}, \frac{3\pi}{2} \right)$
\end{solution}
\end{problem}

\begin{problem}
$\sec^{2} \left( x \right) \leq 4$

\begin{solution}
$\left[ 0, \frac{\pi}{3} \right] \cup \left[ \frac{2\pi}{3}, \frac{4\pi}{3} \right] \cup \left[ \frac{5\pi}{3}, 2\pi \right]$
\end{solution}
\end{problem}

\begin{problem}
 $\cos^{2} \left( t \right) > \frac{1}{2}$

\begin{solution}
$\left[ 0, \frac{\pi}{4} \right) \cup \left( \frac{3\pi}{4}, \frac{5\pi}{4} \right) \cup \left( \frac{7\pi}{4}, 2\pi \right]$
\end{solution}
\end{problem}
 
\begin{problem}
$\cos \left( 2x \right) \leq 0$

\begin{solution}
$\left[ \frac{\pi}{4}, \frac{3\pi}{4} \right] \cup \left[ \frac{5\pi}{4}, \frac{7\pi}{4} \right]$
\end{solution}
\end{problem}

\begin{problem}
$\sin \left( t + \frac{\pi}{3} \right) > \frac{1}{2}$

\begin{solution}
$\left[ 0, \frac{\pi}{2} \right) \cup \left( \frac{11\pi}{6}, 2\pi \right]$
\end{solution}
\end{problem}

\begin{problem}
$\cot^{2} \left( x \right) \geq \frac{1}{3}$

\begin{solution}
$\left( 0, \frac{\pi}{3} \right] \cup \left[ \frac{2\pi}{3}, \pi \right) \cup \left( \pi, \frac{4\pi}{3} \right] \cup \left[ \frac{5\pi}{3}, 2\pi \right)$
\end{solution}
\end{problem}

\begin{problem}
$2\cos(t) \geq 1$ 

\begin{solution}
$\left[0, \frac{\pi}{3}\right] \cup \left[\frac{5\pi}{3}, 2\pi\right]$
\end{solution}
\end{problem}

\begin{problem}
$\sin(5x) \geq 5$ 

\begin{solution}
No solution
\end{solution}
\end{problem}

\begin{problem}
$\cos(3t) \leq 1$

\begin{solution}
$[0, 2\pi]$
\end{solution}
\end{problem}

\begin{problem}
$\sec(x) \leq \sqrt{2}$

\begin{solution}
$\left[0, \frac{\pi}{4} \right] \cup \left(\frac{\pi}{2}, \frac{3\pi}{2}\right) \cup \left[\frac{7\pi}{4}, 2\pi\right]$
\end{solution}
\end{problem}

\begin{problem} \label{firstineqlast}
$\cot(t) \leq 4$ 

\begin{solution}
$\left[\text{arccot}(4), \pi \right) \cup \left[ \pi + \text{arccot}(4), 2\pi\right)$
\end{solution}
\end{problem}

\end{question}

\begin{question}
In Exercises \ref{secondineqefirst} - \ref{secondineqlast}, solve the inequality.  Express the exact answer in \underline{interval} notation, restricting your attention to $-\pi \leq x \leq \pi$.

\begin{problem} \label{secondineqefirst}
$\cos \left( x \right) > \frac{\sqrt{3}}{2}$ 

\begin{solution}
$\left( -\frac{\pi}{6}, \frac{\pi}{6} \right)$
\end{solution}
\end{problem}

\begin{problem}
$\sin(t) > \frac{1}{3}$

\begin{solution}
$\left( \arcsin\left(\frac{1}{3}\right), \pi - \arcsin\left(\frac{1}{3}\right) \right)$
\end{solution}
\end{problem}

\begin{problem} 
$\sec \left( x \right) \leq 2$

\begin{solution}
$\left[ -\pi, -\frac{\pi}{2} \right) \cup \left[ -\frac{\pi}{3}, \frac{\pi}{3} \right] \cup \left( \frac{\pi}{2}, \pi \right]$
\end{solution}
\end{problem}

\begin{problem}
$\sin^{2} \left( t \right) < \frac{3}{4}$

\begin{solution}
$\left( -\frac{2\pi}{3}, -\frac{\pi}{3} \right) \cup \left( \frac{\pi}{3}, \frac{2\pi}{3} \right)$
\end{solution}
\end{problem}

\begin{problem}
$\cot \left( x \right) \geq -1$

\begin{solution}
$\left( -\pi, -\frac{\pi}{4} \right] \cup \left( 0, \frac{3\pi}{4} \right]$
\end{solution}
\end{problem}

\begin{problem} \label{secondineqlast}
$\cos(t) \geq \sin(t)$

\begin{solution}
$\left[ -\frac{3\pi}{4}, \frac{\pi}{4} \right]$
\end{solution}
\end{problem}

\end{question}

\begin{question}
In Exercises \ref{thirdineqfirst} - \ref{thirdineqlast}, solve the inequality.  Express the exact answer in \underline{interval} notation, restricting your attention to $-2\pi \leq x \leq 2\pi$.

\begin{problem} \label{thirdineqfirst}
$\csc \left( x \right) > 1$ 

\begin{solution}
$\left( -2\pi, -\frac{3\pi}{2} \right) \cup \left( -\frac{3\pi}{2}, -\pi \right) \cup \left( 0, \frac{\pi}{2} \right) \cup \left( \frac{\pi}{2}, \pi \right)$
\end{solution}
\end{problem}

\begin{problem}
$\cos(t) \leq \frac{5}{3}$

\begin{solution}
$[-2\pi, 2\pi]$
\end{solution}
\end{problem}

\begin{problem}
$\cot(x) \geq 5$

\begin{solution}
$\left(-2\pi, \text{arccot}(5) - 2\pi\right] \cup \left(-\pi, \text{arccot}(5) - \pi\right] \cup \left(0, \text{arccot}(5)\right] \cup \left(\pi, \pi + \text{arccot}(5)\right]$
\end{solution}
\end{problem}

\begin{problem}
$\tan^{2} \left( t \right) \geq 1$

\begin{solution}
$\left[ -\frac{7\pi}{4}, -\frac{3\pi}{2} \right) \cup \left( -\frac{3\pi}{2}, -\frac{5\pi}{4} \right] \cup \left[ -\frac{3\pi}{4}, -\frac{\pi}{2} \right) \cup \left( -\frac{\pi}{2}, -\frac{\pi}{4} \right] \cup \left[ \frac{\pi}{4}, \frac{\pi}{2} \right) \cup \left( \frac{\pi}{2}, \frac{3\pi}{4} \right] \cup \left[ \frac{5\pi}{4}, \frac{3\pi}{2} \right) \cup \left( \frac{3\pi}{2}, \frac{7\pi}{4} \right]$
\end{solution}
\end{problem}

\begin{problem}
$\sin(2x) \geq \sin(x)$

\begin{solution}
$\left[ -2\pi, -\frac{5\pi}{3} \right] \cup \left[ -\pi, -\frac{\pi}{3} \right] \cup \left[ 0, \frac{\pi}{3} \right] \cup \left[ \pi, \frac{5\pi}{3} \right]$
\end{solution}
\end{problem}

\begin{problem} \label{thirdineqlast}
$\cos(2t) \leq \sin(x)$

\begin{solution}
$\left[ -\frac{11\pi}{6},  -\frac{7\pi}{6} \right] \cup \left[ \frac{\pi}{6}, \frac{5\pi}{6} \right] \cup, \left\{ -\frac{\pi}{2}, \frac{3\pi}{2} \right\}$
\end{solution}
\end{problem}

\end{question}

\begin{question}
In Exercises \ref{invineqfirst} - \ref{invineqlast}, solve the given inequality.

\begin{problem}  \label{invineqfirst} 
$\arcsin(2x) > 0$

\begin{solution}
$\left(0, \frac{1}{2}\right]$
\end{solution}
\end{problem}

\begin{problem}
$3 \arccos(t) \leq \pi$

\begin{solution}
$\left[\frac{1}{2}, 1\right]$
\end{solution}
\end{problem}

\begin{problem}
$6 \, \text{arccot}(7x) \geq \pi$

\begin{solution}
$\left(-\infty, \frac{\sqrt{3}}{7} \right]$
\end{solution}
\end{problem}

\begin{problem}
$\pi > 2\arctan(t)$

\begin{solution}
$(-\infty, \infty)$
\end{solution}
\end{problem}

\begin{problem}
$2\arcsin(x)^2 > \pi \arcsin(x)$

\begin{solution}
$[-1,0)$
\end{solution}
\end{problem}

\begin{problem} \label{invineqlast} 
$12 \arccos(t)^2+2\pi^2>11\pi \arccos(t)$

\begin{solution}
$\left[-1, -\frac{1}{2}\right) \cup \left( \frac{\sqrt{2}}{2}, 1\right]$
\end{solution}
\end{problem}

\end{question}

\begin{question}
In Exercises \ref{domainfirst} - \ref{domainlast}, express the domain of the function using the extended interval notation. (See Example \ref{TrigDomainEx1} and Section \ref{extendedinterval} for details.)

\begin{problem} \label{domainfirst}
$f(x) = \frac{1}{\cos(x) - 1}$

\begin{solution}
$\displaystyle \bigcup_{k=-\infty}^{\infty} \left( 2k\pi, (2k+2)\pi \right)$
\end{solution}
\end{problem}

\begin{problem}
$f(t) = \frac{\cos(t)}{\sin(t) + 1}$

\begin{solution}
$\displaystyle \bigcup_{k=-\infty}^{\infty} \left( \frac{(4k - 1)\pi}{2}, \frac{(4k + 3)\pi}{2} \right)$
\end{solution}
\end{problem}

\begin{problem}
$f(x) = \sqrt{\tan^{2}(x) - 1}$

\begin{solution}
$\displaystyle \bigcup_{k=-\infty}^{\infty} \left\{ \left[ \frac{(4k + 1)\pi}{4}, \frac{(2k + 1)\pi}{2} \right) \cup \left( \frac{(2k + 1)\pi}{2}, \frac{(4k + 3)\pi}{4} \right] \right\}$
\end{solution}
\end{problem}

\begin{problem}
$f(t) = \sqrt{2 - \sec(t)}$

\begin{solution}
$\displaystyle \bigcup_{k=-\infty}^{\infty} \left\{ \left[ \frac{(6k - 1)\pi}{3}, \frac{(6k + 1)\pi}{3} \right] \cup \left( \frac{(4k + 1)\pi}{2}, \frac{(4k + 3)\pi}{2} \right) \right\}$
\end{solution}
\end{problem}

\begin{problem}
$f(x) = \csc(2x)$

\begin{solution}
$\displaystyle \bigcup_{k=-\infty}^{\infty} \left( \frac{k\pi}{2}, \frac{(k+1)\pi}{2} \right)$
\end{solution}
\end{problem}

\begin{problem}
$f(t) = \frac{\sin(t)}{2 + \cos(t)}$

\begin{solution}
$(-\infty, \infty)$
\end{solution}
\end{problem}

\begin{problem}
$f(x) = 3\csc(x) + 4\sec(x)$ 

\begin{solution}
$\displaystyle \bigcup_{k=-\infty}^{\infty} \left( \frac{k\pi}{2}, \frac{(k+1)\pi}{2} \right)$
\end{solution}
\end{problem}

\begin{problem}
$f(t) = \ln\left( |\cos(t)| \right)$

\begin{solution}
$\displaystyle \bigcup_{k=-\infty}^{\infty} \left( \frac{(2k - 1)\pi}{2}, \frac{(2k+1)\pi}{2} \right)$
\end{solution}
\end{problem}

\begin{problem}\label{domainlast}
$f(x) = \arcsin(\tan(x))$ 

\begin{solution}
$\displaystyle \bigcup_{k=-\infty}^{\infty} \left[ \frac{(4k - 1)\pi}{4}, \frac{(4k+1)\pi}{4} \right]$
\end{solution}
\end{problem}

\end{question}

\begin{problem} \label{frequencynumberconnection} 

\begin{enumerate}

\item With the help of your classmates, determine the number of solutions to $\sin(x) = \frac{1}{2}$ in $[0,2\pi)$.  Then find the number of solutions to $\sin(2x) = \frac{1}{2}$,  $\sin(3x) = \frac{1}{2}$ and $\sin(4x) = \frac{1}{2}$ in $[0,2\pi)$. What pattern emerges?   Explain how this pattern would help you solve equations like $\sin(11x) = \frac{1}{2}$.  

\item Repeat the above exercise focusing on  $\sin\left(\frac{x}{2}\right)  = \frac{1}{2}$,  $\sin\left(\frac{3x}{2}\right)  = \frac{1}{2}$ and $\sin\left(\frac{5x}{2}\right)  = \frac{1}{2}$.  What pattern emerges here?  

\item  Replace sine with tangent and $\frac{1}{2}$ with $1$ and repeat the whole exploration.

\end{enumerate}

\end{problem}

\begin{problem}
Suppose an object weighing $10$ pounds is suspended from the ceiling by a spring which stretches $2$ feet to its equilibrium position when the object is attached.  

\begin{enumerate}

\item  Find the spring constant $k$ in $\frac{\text{lbs.}}{\text{ft.}}$ and the mass of the object in slugs.

\begin{solution}
$k = 5 \, \frac{\text{lbs.}}{\text{ft.}}$ and $m = \frac{5}{16} \, \text{slugs}$
\end{solution}

\item  Find the equation of motion of the object if it is released from $1$ foot \textit{below} the equilibrium position from rest.  When is the first time the object passes through the equilibrium position? In which direction is it heading?

\begin{solution}
$x(t) = \sin\left(4t + \frac{\pi}{2}\right)$.  The object first passes through the equilibrium point when $t = \frac{\pi}{8} \approx 0.39$ seconds after the motion starts.  At this time, the object is heading upwards.
\end{solution}

\item  Find the equation of motion of the object if it is released from $6$ inches \textit{above} the equilibrium position with a \textit{downward} velocity of $2$ feet per second.  Find when the object passes through the equilibrium position heading downwards for the third time.

\begin{solution}
$x(t) = \frac{\sqrt{2}}{2} \sin\left(4t + \frac{7\pi}{4}\right)$.  The object passes through the equilibrium point heading downwards for the third time when $t = \frac{17\pi}{16} \approx 3.34$ seconds.
\end{solution}

\end{enumerate}

\end{problem}

\begin{problem} \label{sinecurvesketckex01IPs} 
In Example \ref{sinecurvesketckex01},  $f(x) = 2\sin(x) - \sin(2x)$  restricted to $0 \leq x \leq 2\pi$.  If $f''(x) = 4 \sin(2x) - 2\sin(x)$,  find the inflection points of the graph of $y = f(x)$.

\begin{solution}
The inflection points are:  $\left( \arccos\left(\frac{1}{4}\right), \frac{3 \sqrt{15}}{8} \right)$, $( \pi, 0)$, and $\left( 2\pi - \arccos\left(\frac{1}{4}\right), - \frac{3 \sqrt{15}}{8} \right)$
\end{solution}
\end{problem}

\begin{problem} \label{IPexamplecosineincdec}
In Example \ref{IPexamplecosine},  $f(x) =x-2\cos(x)$ restricted to $0 \leq x \leq 2\pi$.  If $f'(x) = 2\sin(x)+1$, list the open intervals over which $f$ is increasing and decreasing.  Find the local extrema. 

\begin{solution}
$f$ is increasing on $\left(0, \frac{7\pi}{6} \right)$ and again on $\left(\frac{11\pi}{6}, 2\pi \right)$;  $f$ is decreasing on  $\left(\frac{7 \pi}{6}, \frac{11\pi}{6} \right)$; local (absolute) max:  $\left(\frac{7\pi}{6}, \frac{7\pi}{6} + \sqrt{3}\right)$; local min:  $\left(\frac{11\pi}{6}, \frac{11\pi}{6} - \sqrt{3}\right)$
\end{solution}
\end{problem}


\begin{problem} Let $f(x) = e^{-x} \, \sin(x)$ for $x > 0$.

\begin{enumerate}

\item  Use the Squeeze Theorem, Theorem \ref{squeezeth},  to find $\ds{\lim_{x \rightarrow \infty} f(x)}$.  Interpret your answer graphically.

\begin{hint}
Since $-1 \leq \sin (x)  \leq 1$, $-e^{-x} \leq   e^{-x}   \sin (x)   \leq e^{-x}$ \ldots
\end{hint}

\begin{solution}
$\ds{\lim_{x \rightarrow \infty} f(x) = 0}$.  We have a horizontal asymptote $y = 0$. 
\end{solution}

\item  Use the fact that  $f'(x) = e^{-x} \cos(x) - e^{-x} \sin(x)$ to help you find the intervals over which $f$ is increasing and decreasing.

\begin{solution}
\begin{itemize} \item $f$ is increasing on $\left(0, \frac{\pi}{4} \right)$, $\left(\frac{5\pi}{4}, \frac{9 \pi}{4} \right)$, $\left(\frac{13\pi}{4}, \frac{17 \pi}{4} \right)$, \ldots :  

In other words: on $\left(0, \frac{\pi}{4} \right)$ along with $\left( \frac{(8k+5) \pi}{4}, \frac{(8k+9)\pi}{4}\right)$ for $k = 0, 1, 2, 3, \ldots$

\item $f$ is decreasing on$\left(\frac{\pi}{4}, \frac{5 \pi}{4} \right)$, $\left(\frac{9\pi}{4}, \frac{13 \pi}{4} \right)$, $\left(\frac{17\pi}{4}, \frac{21\pi}{4} \right)$, \ldots : 

In other words:  on $\left( \frac{(8k+1) \pi}{4}, \frac{(8k+5)\pi}{4}\right)$ for $k  = 0, 1, 2, 3 \ldots$

\end{itemize}
\end{solution}

\item  Use the fact that $f''(x) = -2 e^{-x} \cos(x)$ to help you find the intervals over which the graph of $f$ is concave up and concave down.

\begin{solution}
\begin{itemize} \item the graph of $f$ is concave up on $\left(\frac{\pi}{2}, \frac{3\pi}{2} \right)$, $\left(\frac{5\pi}{2}, \frac{7 \pi}{2} \right)$, $\left(\frac{9\pi}{2}, \frac{11 \pi}{2} \right)$, \ldots :  

In other words: on  $\left( \frac{(4k+1) \pi}{2}, \frac{(4k+3)\pi}{2}\right)$ for $k = 0,1, 2, 3, \ldots$

\item the graph of $f$ is concave down on  $\left(0, \frac{\pi}{2} \right)$, $\left(\frac{3\pi}{2}, \frac{5 \pi}{2} \right)$, $\left(\frac{7 \pi}{2}, \frac{9\pi}{2} \right)$, \ldots : 

In other words:  on $\left(0, \frac{\pi}{2} \right)$ along with $\left( \frac{(4k+3) \pi}{2}, \frac{(4k+5)\pi}{2}\right)$ for $k  = 0, 1, 2, 3 \ldots$

\end{itemize}
\end{solution}

\end{enumerate}

\end{problem}


\begin{problem} \label{squeezefordampedmotionexercise}
Recall  $x(t) = 10e^{-t/5} \sin\left(t + \frac{\pi}{3}\right)$ from  Example \ref{underdampedresonance} number \ref{underdampedproblem} models underdamped motion. Use the Squeeze Theorem, Theorem \ref{squeezeth}, to prove $\ds{\lim_{t \rightarrow \infty} x(t) = 0}$.

\textbf{HINT:}  Since $-1 \leq \sin\left(t + \frac{\pi}{3}\right) \leq 1$, $-10e^{-t/5} \leq   10e^{-t/5} \sin\left(t + \frac{\pi}{3}\right) \leq 10e^{-t/5}$ \ldots

\end{problem}

\end{document}


